\chapter{The theorems}

Two theorem families anchor the project: the Cender attention theorems (design theorems for the
constant-memory attention line) and the TransCender wire-identity theorems (the checkpoint
identity contract). A design theorem is not a measured result; the distinction is kept.

\section{Theorem 1 --- constant memory}

\textbf{Statement.} For any sequence length $n$, Cender state memory
$|\text{state}| = d^2 \cdot \text{sizeof(float)}$ is constant. Transformers KV memory
$|\text{KV}| = 2nd\cdot\text{sizeof(float)}$ is linear in $n$.

\begin{quote}
\textbf{Proof.} The Cender recurrence updates a fixed matrix in place:
$\text{state}_{t+1} = \gamma\cdot\text{state}_t + k_t \otimes e_t$ where
$k_t = W_k x_t$, $e_t = v_t - \text{state}_t k_t$, $\gamma \in [0,1]$. No term depends on $n$;
the allocation is one \texttt{Tensor\{d,d\}}, never appended per token. The baseline appends
$K_t, V_t$ per token, so its state grows as $\Theta(n)$. At $d{=}4096$, $n{=}4$M:
$|\text{KV}| \approx 2\cdot4\text{M}\cdot4096\cdot2\,\text{B} \approx 64$~GB (FP16) versus
$|\text{state}| = 4096^2 \cdot 4\,\text{B} \approx 64$~MB --- a ratio of ${\approx}1000\times$,
constant in $n$. \hfill$\square$
\end{quote}

\textbf{Verification gate.} \texttt{bench\_inference} tok/s vs $n$ flat for Cender, decreasing
for the baseline; max RSS flat within 5\% at $n \in \{4K, 32K, 128K, 1M, 4M\}$; no PSNR
regression.

\section{Theorem 2 --- no context rot by construction}

\textbf{Statement.} Under Cender, recall accuracy $\mathrm{Acc}(pos)$ is independent of total
length $n$ and needle position $pos$, up to decay tuning. Under softmax attention,
$\mathrm{Acc}(pos)$ is U-shaped and degrades with $n$.

\textbf{Proof sketch.}
\begin{enumerate}
\item \textbf{Softmax dilution lemma.} $\mathrm{softmax}(QK^T/\sqrt{d})$ distributes mass with
$\sum_i p_i = 1$. As $n\to\infty$, expected $p_{\text{needle}} \to 1/n \to 0$; cross-entropy on
the middle token scales as $-\log p_{\text{needle}} \propto \log n$. Middle recall collapses ---
the context-rot family.
\item \textbf{Delta associative lemma.} \texttt{state} is an associative memory over key-value
bindings via the delta rule, with recall error $O(1/d)$ independent of $n$ (DeltaNet Thm.~3.1,
arXiv 2406.06484; Titans Lemma~2). Retrieval $y = \text{state}\cdot q$ is one matrix-vector
product --- no softmax over $n$. Orthogonal page writes add bindings in an orthogonal subspace,
so interference $\langle \text{old}, \text{new}\rangle \approx 0$.
\item \textbf{Decay controls forgetting, not dilution.} $\gamma^t$ fades stale state
geometrically, preserving recent context without mass dilution. Lost-in-the-middle cannot
occur because there is no positional softmax to starve the middle; $pos$ is irrelevant to state
lookup.
\item \textbf{Construction.} There exists $\gamma$ and $k_{\text{orth}}$ such that for any $n$
up to 1T tokens, $\mathrm{Acc}_{\text{Cender}}(n, pos) \ge
\mathrm{Acc}_{\text{TF}}(4K, pos) - \varepsilon$ with $\varepsilon \to 0$. \hfill$\square$
\end{enumerate}

\textbf{Verification gates.} Needle test: 100-page PDF, secret on page 47, positions swept at
0/25/50/75/100\%; Cender flat $\ge 95\%$ at ${\sim}100$K tokens with 95\% garbage. Context-rot
test: full dump vs the RLM router filter; 70\% cost cut, 10$\times$ latency target.

\section{Theorem 3 --- bounded reads, fail-closed parsing}

Every \texttt{.idx} reader bounds-checks the header against the magic lengths (22 B for
\texttt{Transformers-FusionIDX}, 15 B for \texttt{TransCenderIDX}) before reading further; a
truncated header fails closed. Proof: the per-magic size gates (30/23) plus a header-bounds
guard, fixed in audit 1; the regression is \texttt{test\_sha256\_corrupt} (23/0 pass).
\hfill$\square$

\section{Theorem 4 --- three-source identity agreement}

The engine registry, the model manifest and the decision log agree on one identity: the engine
accepts \texttt{Transformers-FusionIDX} (22 B) plus legacy \texttt{TransCenderIDX} (15 B);
\texttt{SolIDX} is model-manifest-only. Proof: D-30 states it, the registry implements it,
Test 8 (28 tests, 0 failures) verifies it. \hfill$\square$

\section{The GPTQ-on-loops obligation}

\textbf{Statement.} On a looped model, a PTQ method that calibrates on step-0 activations only
is dominated by Hessian accumulation across recurrence steps.

\begin{quote}
\textbf{Proof sketch.} One-step GPTQ's Hessian $H = E[a_0a_0^T]$ weights only the step-0 input
directions. The loop visits directions $a_t$ ($t>0$) nearly orthogonal to $a_0$; their
quantization error feeds back through the recurrence (feedback exposure). Accumulating
$H = \sum_t E[a_t a_t^T]$ weights every visited direction; the published result: one-step GPTQ
loses to RTN on 5/9 looped checkpoints while accumulation wins 9/9 (2609.30820).
\hfill$\square$
\end{quote}
